\section{Erd\H{o}s Problem 72: a sparse unavoidable set of cycle lengths}

\subsection{1) Statement and construction}
Does there exist a set $A\subseteq\mathbb N$ of density zero and a
constant $c$ such that every sufficiently large finite graph of
average degree at least $c$ has a cycle with length in $A$?
Take the explicit set
\[
 A=\{2^j:j\geq2\}.
\]
The assertion for this set follows from a theorem of Liu and
Montgomery, strengthening Verstraete's earlier existence result.
We use that theorem explicitly and prove the numerical deduction.

\subsection{2) External theorem and the required uniform constant}
Theorem 1.1 of Liu--Montgomery gives $d_0$ such that every graph
of average degree $d\geq d_0$ has a number
\[
 \ell\geq\frac{d}{10(\log d)^{12}}
\]
with cycles of every even length in $[(\log\ell)^8,\ell]$.
Choose $L\geq4$ so large that $(\log x)^8\leq x/2$ for all
$x\geq L$. Then choose $c\geq d_0$ so large that
$d/[10(\log d)^{12}]\geq L$ whenever $d\geq c$.
Both choices are possible by the elementary comparison between powers
of logarithm and positive powers of $x$.

Let $G$ have average degree at least $c$, and take the corresponding
$\ell$. If $j=\lfloor\log_2\ell\rfloor$, then
\[
 \ell/2<2^j\leq\ell,
 \qquad 2^j\geq4,
 \qquad 2^j\geq(\log\ell)^8.
\]
Thus $2^j$ is one of the even cycle lengths supplied by the theorem.
It belongs to $A$, as required. In fact this argument works for
every finite graph meeting the average-degree threshold, without
an additional lower bound on its order.

\subsection{3) Density and a useful extension}
For $N\geq4$,
\[
 |A\cap[1,N]|=\lfloor\log_2N\rfloor-1,
 \qquad \frac{|A\cap[1,N]|}{N}\longrightarrow0.
\]
Hence the chosen set really has natural density zero; no appeal to
its complement being large is needed.

More generally, let $a_1<a_2<\cdots$ be even integers with
$a_{j+1}\leq K a_j$ for some fixed $K>1$. For sufficiently large
$\ell$, take the largest $a_j\leq\ell$. Then
$\ell<a_{j+1}\leq K a_j$, so $a_j>\ell/K$. Once
$(\log\ell)^8\leq\ell/K$ and $\ell\geq a_1$, the same interval
theorem forces a cycle length from this sequence. This explains why
exponential sparsity is compatible with being unavoidable.

\textbf{Outcome: SOLVED, known affirmative resolution reconstructed
using the stated external theorem.} The theorem's graph-theoretic
proof is not supplied or independently formally verified here.

\subsection{4) Sources}
\begin{itemize}
\item H. Liu and R. Montgomery, \emph{A solution to Erd\H{o}s and
Hajnal's odd cycle problem}, Theorem 1.1,
\texttt{https://arxiv.org/abs/2010.15802}.
\item T. F. Bloom, Problem 72,
\texttt{https://www.erdosproblems.com/72}, checked 23 September 2026.
\end{itemize}
The estimate measures coverage using the explicitly cited theorem,
not novelty of the construction.
\subsection{5) COMPLETION ESTIMATE}
\textbf{COMPLETION: 100\%}
